\chapter{Il nucleo del control plane}
\label{cap:06-nucleo}

Il nucleo \`e ci\`o che resta quando si compila \nf{} senza alcun modulo di coda:
i quattro obblighi di una piattaforma FaaS enunciati nel
\cref{cap:02-stato-dell-arte} vi sono tutti implementati, e il modulo pu\`o solo
aggiungere politiche. Il codice non sta pi\`u in un solo progetto: la macchina a
stati dei tentativi, lo store delle esecuzioni e il registro della capacit\`a
vivono in \code{platform/libs/execution-runtime}, i contratti verso i moduli in
\code{platform/libs/control-plane-spi}, e il resto (registro delle funzioni, API,
dispatcher, metriche) in \code{platform/control-plane}
(\cref{cap:08-sistema-moduli}). Insieme alle tre librerie obbligatorie il modulo
\code{control-plane} conta 17\,334 righe non vuote di produzione, contro le 6\,299 dei
dieci moduli opzionali (misura del 28~settembre~2026, la stessa del
\cref{cap:01-introduzione}). Il capitolo attraversa i componenti
nell'ordine in cui li incontra una richiesta, soffermandosi sulle decisioni che il
codice motiva esplicitamente. Le garanzie sul ciclo di vita, che attraversano pi\`u
componenti, sono raccolte in un contratto scritto
(\code{docs/architecture/0001-execution-lifecycle-contract.md}, nel seguito ADR~0001)
che il capitolo cita per nome.

\section{Registro e risoluzione della specifica}

\subsection{\code{FunctionRegistry}, \code{FunctionService} e il catalogo persistente}

Il registro non \`e pi\`u una mappa mutabile: \code{FunctionRegistry} conserva uno
\emph{snapshot} immutabile, sostituito per intero a ogni scrittura, con due viste
--- quella pubblica, che le API leggono, e quella di recupero, che include anche
le funzioni indisponibili (in rimozione parziale o non ancora ripristinate).
La lettura non prende lock; le scritture sono \code{synchronized} e pubblicano
lo snapshot solo dopo averlo reso durevole.

La durabilit\`a \`e la novit\`a pi\`u rilevante rispetto alla prima versione del
progetto. \code{FunctionCatalog} scrive l'intero catalogo in un file
(\code{nanofaas.registry.path}): file temporaneo nella stessa directory, permessi
\code{0600} sul file e \code{0700} sulla directory (in modo best-effort, perch\'e un
volume montato da root non lo consente), \code{fsync} del file, spostamento atomico
e \code{fsync} della directory. Il commento nel sorgente nomina il costo senza
nasconderlo: ogni mutazione riscrive e sincronizza tutto il catalogo, serializzata
sul monitor del registro, e la scrittura differita con flush coalescente \`e la
via di miglioramento se la frequenza diventa un collo di bottiglia. Un catalogo
illeggibile o con schema sconosciuto fa fallire l'avvio, anzich\'e ripartire vuoto.
Nel chart Helm il percorso sta su un volume persistente (\code{/var/lib/nanofaas}).

Sopra il registro, \code{FunctionService} orchestra le operazioni che la
registrazione comporta oltre alla memorizzazione, tutte sotto un lock per nome
(\code{FunctionOperationLocks}) e tutte rifiutate finch\'e il catalogo non \`e
stato ripristinato (\code{FunctionRestoreGate}). L'ordine della registrazione
\`e esplicito:

\begin{enumerate}
  \item risoluzione dei default (\cref{sec:06-resolver});
  \item rifiuto se il nome \`e in rimozione parziale o gi\`a presente;
  \item validazione dell'immagine attraverso il validatore del backend attivo;
  \item provisioning del deployment gestito, se la modalit\`a lo prevede;
  \item notifica ai listener di registrazione (capacit\`a, code, metriche);
  \item commit durevole nel registro --- l'ultimo passo.
\end{enumerate}

Se un passo fallisce, i listener gi\`a notificati vengono ritirati e le risorse
gi\`a create vengono rilasciate: la funzione non compare a met\`a. La cancellazione
percorre l'ordine inverso e mette il commit durevole \emph{per ultimo}.

La registrazione non \`e idempotente: registrare un nome gi\`a presente
restituisce \code{409 Conflict}. La modifica avviene invece per
\code{PATCH}, con una regola dichiarata nel tipo della richiesta:

\begin{lstlisting}[language=Java, caption={\code{FunctionUpdateRequest}: il
commento nel sorgente enuncia la regola.}]
/**
 * Partial update of a registered function. Only fields that the control plane can apply without
 * touching the deployment are accepted; a null field leaves the current value alone.
 *
 * <p>Unknown properties are rejected rather than ignored, so patching an immutable field such as
 * {@code image} or {@code executionMode} fails loudly instead of appearing to succeed.</p>
 */
@JsonIgnoreProperties(ignoreUnknown = false)
public record FunctionUpdateRequest(
        @Min(1) Integer concurrency, @Min(1) Integer timeoutMs,
        @Min(0) Integer maxRetries, ConcurrencyControlConfig concurrencyControl
) { ... }
\end{lstlisting}

La superficie modificabile a caldo \`e quindi esattamente quattro campi, e sono
i quattro che governano il comportamento del control plane senza toccare il
deployment sottostante. Il rifiuto delle propriet\`a sconosciute \`e la parte
importante: senza di esso, un \code{PATCH} che tentasse di cambiare
\code{image} riceverebbe \code{200} e non cambierebbe nulla, che \`e il modo
peggiore di fallire. Per una funzione gestita, l'aggiornamento raggiunge anche il
provider: il proxy del backend container deriva dalla specifica il timeout del
singolo salto e il limite di ammissione, e senza questo passo il deployment
continuerebbe a imporre i valori originali mentre il chiamante crede in quelli
modificati. Un errore di un listener dopo l'aggiornamento \emph{non} produce un
rollback (il commento lo giustifica: il percorso di rollback della registrazione
cancella la coda della funzione, che sarebbe peggio di un listener che manca un
aggiornamento; i listener sono idempotenti, quindi ripetere lo stesso \code{PATCH}
converge).

Una nota sulla semantica del blocco annidato: \code{concurrencyControl} viene
sostituito \emph{in blocco}, non fuso campo per campo. La ragione \`e che i campi
di quel blocco hanno significati diversi a seconda del modo
(\cref{cap:12-concurrency-control}), e una fusione parziale produrrebbe
configurazioni ibride prive di senso --- per esempio una soglia di carico
ereditata da \code{ADAPTIVE\_PER\_POD} sopravvissuta al passaggio a
\code{BUDGETED}.

\paragraph{Rimozione parziale.} Una cancellazione pu\`o finire in tre stati, e solo
due possono essere presentati come esito affidabile. La rimozione completata
cancella la voce dal catalogo per ultima. Se il backend segnala una
\code{PartialDeprovisionException} --- risorse che non ha potuto liberare --- la
funzione entra in \emph{rimozione in sospeso}: la voce resta nel catalogo (\`e ci\`o
che tiene tracciabili i residui e fa s\`i che un \code{DELETE} ripetuto riprenda la
pulizia invece di rispondere \code{404}), i listener \emph{non} vengono
riattivati, e ogni via d'ingresso alla funzione --- invocazione, enqueue, patch,
repliche, ri-registrazione --- risponde \code{409} con codice
\code{FUNCTION\_REMOVAL\_PENDING}. Il rollback operativo \`e ammesso solo quando
nulla di necessario \`e andato perso. La rimozione in sospeso \`e uno stato di
processo, non persistito: dopo un riavvio la funzione viene ripristinata come
qualunque altra, e se la pulizia \`e ancora voluta il \code{DELETE} va ripetuto
(ADR~0001, \S 8.3). \`E una conseguenza del vincolo di un singolo pod con stato in
memoria, non una svista.

\subsection{\code{FunctionSpecResolver}: dove i default diventano valori}
\label{sec:06-resolver}

La risoluzione trasforma una specifica parzialmente compilata in una completa.
\`E il punto in cui il modello ``\code{null} significa non specificato''
(\cref{cap:05-contratti-dati}) viene consumato:

\begin{lstlisting}[language=Java]
Optional.ofNullable(spec.timeoutMs()).orElse(defaults.timeoutMs()),
Optional.ofNullable(spec.concurrency()).orElse(defaults.concurrency()),
Optional.ofNullable(spec.queueSize()).orElse(defaults.queueSize()),
Optional.ofNullable(spec.maxRetries()).orElse(defaults.maxRetries()),
\end{lstlisting}

I default delle politiche di concorrenza non vivono pi\`u nel resolver: il resolver
delega a \code{ConcurrencyControlConfig.normalize}, nel progetto dei contratti
condivisi, cos\`i che la stessa normalizzazione valga ovunque il blocco venga letto.
Due default sono commentati in modo che vale la pena riportare perch\'e esprimono
principi, non valori.

\paragraph{Le soglie del controllore adattivo.}
\begin{lstlisting}[language=Java]
// Fractions of latency degradation over the function's best observed service time:
// back off above 2x, grow again below ~1.18x. See the concurrency-control module.
private static final double DEFAULT_HIGH_LOAD_THRESHOLD = 0.5;
private static final double DEFAULT_LOW_LOAD_THRESHOLD = 0.15;
\end{lstlisting}

Le soglie sono espresse come frazioni di degradazione rispetto al \emph{miglior}
tempo di servizio osservato per quella funzione, non come valori assoluti. \`E
la stessa costruzione di TCP Vegas~\cite{brakmo1995vegas}: il riferimento \`e il
minimo storico, assunto come stima del costo di servizio irriducibile, e il
segnale \`e lo scostamento da esso. Il \cref{cap:12-concurrency-control} discute
perch\'e questa costruzione, corretta nel contesto originale, si riveli
insufficiente qui.

\paragraph{L'SLO di default.}
\begin{lstlisting}[language=Java]
// A generic service-time SLO for a function that did not state one. Deliberately not
// derived from anything the platform measures: a target the controller inferred from
// observed latency would move whenever the function got slower, which is the one thing
// an SLO must not do.
private static final long DEFAULT_TARGET_LATENCY_MS = 250L;
\end{lstlisting}

L'osservazione \`e sottile e generale: un obiettivo derivato dall'osservazione
non \`e un obiettivo. Se il sistema stimasse l'SLO dalla latenza misurata, un
degrado prestazionale sposterebbe automaticamente l'asticella, e il controllore
si dichiarerebbe conforme mentre peggiora. Il default \`e quindi un numero
arbitrario e dichiarato tale.

\paragraph{Default dipendenti dalla modalit\`a.} La risoluzione del blocco di
scaling avviene solo per le funzioni in modalit\`a \code{DEPLOYMENT}: per
\code{LOCAL} e \code{EXTERNAL} il blocco resta come fornito, perch\'e il control
plane non gestisce repliche che non ha creato. Per \code{DEPLOYMENT} senza blocco
esplicito il default \`e strategia \code{INTERNAL}, da una a dieci repliche,
metrica \code{queue\_depth} con obiettivo 5. Sono validati anche i limiti
(\code{minReplicas} non negativo, \code{maxReplicas} almeno 1, minimo non
superiore al massimo).

Le metriche ammesse per la strategia interna sono validate contro un insieme
chiuso (\code{queue\_depth}, \code{in\_flight}, \code{rps}): una metrica
sconosciuta fa fallire la registrazione con \code{400} anzich\'e produrre un
autoscaler che non scala e non lo dice.

\section{Capacit\`a di ammissione}
\label{sec:06-capacita}

Prima di creare un'esecuzione il nucleo riserva le risorse che essa
tratterr\`a. Le riserve sono finite per configurazione, e i valori di default
sono calibrati sul contenitore minimo supportato dal chart (512~MiB):
\code{nanofaas.invocation-capacity.*}.

\begin{center}
\small
\begin{adjustbox}{max width=\textwidth}
\begin{tabular}{@{}lll@{}}
\toprule
\textbf{Risorsa} & \textbf{Globale / per funzione} & \textbf{Che cosa conta} \\
\midrule
esecuzioni & 4\,096 / 512 & proprietari logici di un'esecuzione viva \\
input canonico & 128~MiB / 32~MiB & rappresentazione Java trattenuta \\
copie fisiche & 64~MiB / 16~MiB & materializzazioni di runtime e trasporto \\
attese (waiter) & 8\,192 / 1\,024 & chiamanti sincroni attaccati, repliche incluse \\
corpo in ingresso & 1~MiB (unico) & byte ricevuti prima del parsing \\
\bottomrule
\end{tabular}
\end{adjustbox}
\end{center}

L'ordine sul percorso di una invocazione HTTP \`e fissato: confine del corpo
(\code{413}), limitatore di frequenza, decodifica JSON e validazione, risoluzione
della funzione, e solo poi la ricerca della chiave di idempotenza. Una \code{Content-Length}
dichiarata oltre il limite \`e rifiutata senza nemmeno sottoscrivere il corpo; per
un corpo a chunk, il primo buffer che supera il confine viene rilasciato e la
sottoscrizione cancellata. Il vincolo che ne segue \`e importante: una
\emph{replica} non pu\`o aggirare il limite del corpo, perch\'e il controllo
precede la ricerca. Una richiesta genuinamente nuova reclama la chiave,
canonicalizza l'input, riserva i proprietari e pubblica il record vivo \emph{prima}
che un waiter sincrono sia ammesso; se l'ammissione del waiter fallisce,
record, chiave e proprietari sono annullati prima del \code{429}.

Una saturazione di esecuzioni, input, copie o waiter risponde \code{429} con
\code{Retry-After: 1} e un corpo stabile,
\code{\{"error":"invocation\_quota\_exceeded","resource":"execution|input|input\_copy|waiter"\}}.
Un \code{413} non \`e sovraccarico. Le stime coprono gli oggetti Java trattenuti dal
control plane, non la memoria nativa dei buffer Netty n\'e quella remota della
funzione: sono una busta di ammissione, non una garanzia sull'RSS del processo.
Quando il modulo \code{runtime-config} \`e attivo, le quattro coppie
globale/per funzione sono modificabili a caldo (\cref{cap:14-runtime-config}); abbassare
un tetto sotto l'occupazione attuale non espelle nulla, e nuove riserve sono
rifiutate finch\'e il drenaggio non riporta l'uso entro il nuovo limite.

\section{Il ciclo di vita di un'esecuzione}

\subsection{\code{ExecutionRecord}}

\code{ExecutionRecord} \`e l'unica struttura mutabile significativa del nucleo.
Contiene lo stato, i timestamp, l'esito, i flag di cold start, la busta di
risposta, i proprietari delle risorse d'input e --- centralmente --- il
\code{CompletableFuture} su cui i chiamanti sincroni attendono.

La mutabilit\`a \`e governata da invarianti dichiarati nel codice.

Il primo \`e che ogni transizione avviene sotto il monitor dell'oggetto e che gli
stati terminali sono definitivi:

\begin{lstlisting}[language=Java]
/**
 * Terminal states (SUCCESS, ERROR, TIMEOUT) are final; every transition between
 * non-terminal states (including RUNNING -> QUEUED for retries) is allowed.
 */
\end{lstlisting}

La permissivit\`a sulle transizioni non terminali \`e ci\`o che rende possibile
la riprova: un'esecuzione fallita torna in coda e ripercorre il ciclo. La
rigidit\`a su quelle terminali \`e ci\`o che rende sicuro il completamento
concorrente: se il dispatcher e la callback arrivano entrambi --- il che accade
quando la funzione risponde \emph{e} invia la callback --- il secondo viene
scartato invece di sovrascrivere l'esito.

Il secondo invariante \`e che il futuro non viene mai completato tenendo il lock.
Oggi il record \emph{sceglie} la risposta terminale sotto il proprio monitor e la
\emph{pubblica} altrove, con un metodo dedicato:

\begin{lstlisting}[language=Java]
/** Publish only the selected terminal answer; callbacks run outside the monitor. */
public void publishTerminal() {
    ... // sotto il monitor: sceglie risultato o eccezione
    if (failure != null) completion.completeExceptionally(failure);
    else completion.complete(result);
}
\end{lstlisting}

Completare un \code{CompletableFuture} esegue in linea le continuazioni dei
sottoscrittori. Farlo sotto il monitor del record significherebbe eseguire codice
applicativo arbitrario --- la mappatura della risposta, la scrittura sulla socket
--- mentre si tiene un lock su cui il completamento concorrente contende.
Separare la scelta dalla pubblicazione ha un secondo effetto: due completamenti
concorrenti scelgono la stessa risposta, perch\'e la scelta \`e una sola.

Due valori sono catturati \emph{una volta} alla costruzione e non riletti dal
task corrente: se l'esecuzione sar\`a leggibile dopo la fine
(\code{readableAfterFinishing}: vero per le asincrone e per quelle con chiave, falso
per una sincrona semplice, la cui risposta \`e gi\`a tornata sulla connessione che il
chiamante teneva) e la chiave di idempotenza. La ragione \`e la stessa per
entrambi: la riprova sostituisce il task con uno privo di chiave, e rileggere il
task dopo di essa declasserebbe proprio le esecuzioni che hanno avuto
problemi, facendo smettere silenziosamente di funzionare l'idempotenza per esse.

\subsection{Il proprietario della transizione terminale}

Fino a un certo punto la transizione terminale era divisa fra pi\`u collaboratori:
lo store archiviava l'esito, e lo spostamento della chiave di idempotenza verso il
suo stato terminale era un listener in una catena best-effort. La conseguenza, nel
sorgente, \`e chiamata per nome (\emph{review finding R2}): un esito troppo grande,
espulso fra i due passi, lasciava una chiave ancora \code{published} che una
replica poteva rivendicare, autorizzando una seconda esecuzione di una funzione
gi\`a eseguita.

\code{ExecutionLifecycle.settle} \`e oggi l'unico proprietario della transizione, e
i passi hanno un ordine che nessun listener pu\`o saltare:

\begin{enumerate}
  \item \emph{protezione della chiave}: la chiave diventa non rivendicabile
    (\code{markTerminal}) prima che l'ultimo riferimento vivo possa sparire, sotto
    il monitor del record, lo stesso che il percorso di ammissione usa per
    pubblicare, cos\`i che un record terminale non possa mai essere pubblicato sotto
    una chiave rivendicabile;
  \item \emph{pubblicazione} della risposta terminale scelta dal record;
  \item \emph{archiviazione} dell'esito, se il suo peso rientra nel budget, e
    rimozione del record vivo;
  \item \emph{notifica} degli osservatori best-effort (metriche, conclusione
    end-to-end): un listener che lancia un'eccezione non pu\`o interrompere la
    pulizia n\'e impedire a un altro proprietario di chiudere.
\end{enumerate}

Il metodo \`e idempotente (un record non terminale \`e ignorato, uno gi\`a assestato
non rifa nulla). Non c'\`e un lock globale sul percorso caldo: il record \`e il proprio
monitor, e nessun listener, dispatcher o codice esterno gira mentre lo si tiene.
\code{ExecutionStore.settle} resta come adattatore per i chiamanti dei moduli e
fallisce esplicitamente se un record \emph{con chiave} viene assestato senza un
proprietario collegato, anzich\'e perdere in silenzio la garanzia di deduplicazione.

\subsection{Il contratto del ciclo di vita}

Le regole che attraversano pi\`u componenti sono state raccolte nell'ADR~0001, con
dieci invarianti (I1--I10). Quattro meritano di essere riportate perch\'e
cambiano cosa un lettore pu\`o assumere.

\paragraph{Il timeout di un chiamante non \`e il timeout dell'esecuzione.}
Un waiter che scade (\code{X-Timeout-Ms}, altrimenti \code{spec.timeoutMs}) ottiene
la propria risposta \code{408}, ma il record condiviso, la chiave, lo store, il
lease e i contatori non cambiano: l'esecuzione prosegue, e una replica successiva
o un polling ne osserva l'esito reale. Il codice lo rende esplicito
sottoscrivendo il futuro condiviso con \code{suppressCancel}, cos\`i che n\'e il
timeout n\'e la disconnessione di un singolo chiamante possano cancellare il lavoro
condiviso. Il \code{TIMEOUT} condiviso \`e riservato alle scadenze di livello
esecuzione, come l'attesa massima in coda (\cref{cap:10-sync-queue}).

\paragraph{Quattro orologi.}
\begin{center}
\small
\begin{adjustbox}{max width=\textwidth}
\begin{tabular}{@{}lll@{}}
\toprule
\textbf{Orologio} & \textbf{Chi lo possiede} & \textbf{Effetto alla scadenza} \\
\midrule
budget del waiter & il chiamante & solo quel waiter riceve \code{408} \\
scadenza del tentativo & il dispatch (\code{timeoutMs}) & segue la politica di riprova \\
politica di riprova & l'esecuzione & nuovo tentativo o errore terminale \\
\code{maxLifetime} & lo store & \code{EXECUTION\_EXPIRED} fabbricato \\
\bottomrule
\end{tabular}
\end{adjustbox}
\end{center}

\paragraph{Tentativo, lease e generazione.} Un tentativo \`e la singola
spedizione di un'esecuzione, e possiede un \code{DispatchLease} sotto la
\emph{generazione} corrente della funzione. Il rilascio del lease \`e idempotente
ed \`e guidato dal futuro \emph{grezzo} del trasporto, non dallo specchio della
scadenza: un handler LOCAL non interrompibile che continua a girare dopo la
scadenza continua a occupare la capacit\`a che gli \`e attribuita, e questa non
viene ceduta a un secondo tentativo. Il futuro terminale e la fine fisica del
lavoro sono due momenti diversi e solo il primo \`e una transizione di stato.
Una funzione rimossa e ri-registrata con lo stesso nome ha una nuova
generazione: un lease vecchio rilascia solo lo slot della vecchia, e una
callback tardiva non pu\`o ricreare stato n\'e toccare le metriche della nuova
(\code{ACTIVE}$\to$\code{RETIRING}$\to$\code{CLOSED}; ADR~0001 \S 8.2).

\paragraph{La cancellazione richiesta prima che esista l'handle.} Una scadenza
amministrativa o uno shutdown possono correre contro un dispatch che non ha
ancora restituito il proprio futuro. La richiesta viene ricordata sul record e
applicata quando l'handle arriva; non \`e mai scartata per essere arrivata presto.

\subsection{\code{ExecutionStore}: vivi e postumi}

La prima versione descritta in questo capitolo era una mappa concorrente con un
thread janitor che ogni minuto eliminava le voci scadute, con tre orizzonti
(\code{ttl}, \code{cleanupTtl}, \code{maxLifetime}). Il sorgente attuale racconta
perch\'e \`e stata sostituita, con due misure nel commento della classe: in una
esecuzione a 1\,093 ammissioni al secondo, \code{function\_inFlight} ha toccato un
massimo di \textbf{2} mentre lo store ne conteneva \textbf{165\,786}: il 99,999\%
di ci\`o che conservava era postumo, e si portava dietro l'apparato dei vivi (il
futuro, il task, la richiesta, l'insieme dei tentativi rilasciati): 35 oggetti
per record, cinque milioni e ottocentomila in tutto, ciascuno da percorrere a
ogni raccolta completa. Il janitor impiegava 12,4~ms su 166\,000 record. E la
ritenzione era dichiarata in tempo e illimitata in numero, quindi la memoria
cresceva con il tasso d'arrivo: il commento riporta 1,05~GB, il 50,6\% del tempo
in GC e una probe di liveness mancata tre volte di seguito.

Le due vite sono oggi due strutture, entrambe cache Caffeine.

\begin{center}
\small
\begin{adjustbox}{max width=\textwidth}
\begin{tabular}{@{}lll@{}}
\toprule
\textbf{Struttura} & \textbf{Contiene} & \textbf{Scadenza e limite} \\
\midrule
\code{inFlight} & \code{ExecutionRecord} mutabili &
  \code{maxLifetime} (30~min) dalla scrittura \\
\code{outcomes} & \code{Outcome} piatto e immutabile &
  \code{ttl} (5~min) se leggibile, \code{sync-ttl} (30~s) altrimenti; \\
 & & tetto in \emph{byte} (\code{max-outcome-bytes}) \\
\bottomrule
\end{tabular}
\end{adjustbox}
\end{center}

Un \code{Outcome} \`e quattro oggetti e 116 byte (contro 35 oggetti e 1\,330 byte di
un record terminale, misurati su 200\,000 record il 2026-08-26 secondo il commento),
ed \`e esattamente ci\`o che serve a due scopi: rispondere a
\code{GET /v1/executions/\{id\}} e servire la replica di una chiave. I tempi sono
\code{long} e non \code{Instant} per la stessa ragione di peso.

Il limite \`e in byte e non in numero, perch\'e un esito leggibile trattiene il
payload del chiamante: 20\,000 esiti da 64~KB occupano 1,28~GB (misura citata nel
sorgente da \code{docs/experiments/control-plane-tuning-2026-09/RESULTS.md}).
Il peso \`e una \emph{stima} conservativa calcolata una volta all'inserimento da
\code{OutcomeWeigher}, che congela in profondit\`a il payload in una copia
immutabile e non condivisa (cos\`i che nulla possa mutarlo, e farne crescere la
dimensione, dopo che \`e stato pesato) e non lo riserializza mai. La visita \`e
limitata in profondit\`a e ampiezza, ma \emph{mai} valuta zero la parte saltata:
se la dimensione non pu\`o essere limitata --- struttura troppo profonda,
contenitore troppo largo, ciclo, riferimento condiviso, oggetto opaco --- l'esito
non viene ritenuto. Anche un esito che da solo supera l'intero budget \`e
rifiutato, anzich\'e inserito per essere espulso un istante dopo. Il default di
\code{max-outcome-bytes} \`e derivato: \code{max-outcomes} $\times$ 116 byte, cio\`e
il budget \`e il vincolo che davvero limita lo heap.

Il rifiuto di ritenere il payload non \`e la perdita della deduplicazione: il
tombstone \`e la chiave, non il payload, ed \`e il proprietario della transizione
a proteggerlo prima di tutto. Una replica di una chiave il cui esito \`e stato
espulso per capacit\`a riceve \code{410 Gone} (\code{OutcomeGoneException}); la
funzione \emph{non} viene rieseguita.

La scadenza amministrativa (\code{maxLifetime}) \`e il secondo orizzonte che
resta, e la sua condizione \`e dichiarata: un'esecuzione la cui callback non
arriva mai non raggiunge uno stato terminale, quindi gli orizzonti ancorati
all'istante di completamento non la toccherebbero mai. La cache dei vivi ha uno
scheduler proprio, con la motivazione riportata nel commento: senza di esso
Caffeine controlla la scadenza solo quando qualcosa tocca la cache, e un record
bloccato su un dispatch perso, che nessuno rilegge pi\`u, restava scaduto ma vivo
a tempo indeterminato --- e con esso lo slot di concorrenza che teneva.
L'espulsione notifica i listener di scadenza amministrativa; il coordinatore dei
tentativi fabbrica un esito \code{EXECUTION\_EXPIRED}, marcato come tale per
distinguerlo da un errore genuino del runtime, e lo assesta.

Due metriche misurano la distanza fra ci\`o che la piattaforma ricorda e ci\`o che
sta eseguendo: \code{execution\_store\_size} (esiti archiviati) e
\code{execution\_in\_flight\_records}. Senza la seconda, la prima potrebbe essere
scambiata per lavoro in corso.

\section{Idempotenza}

L'idempotenza \`e attivata dall'header \code{Idempotency-Key} ed \`e implementata
da \code{IdempotencyStore}, una cache Caffeine. La chiave \`e composta con il nome
della funzione: due funzioni con lo stesso valore di chiave sono associazioni
indipendenti. Il problema che risolve non \`e la deduplicazione a posteriori ma la
\emph{corsa}: due richieste con la stessa chiave che arrivano contemporaneamente
devono produrre una sola esecuzione.

Il meccanismo \`e una rivendicazione a fasi, e le fasi sono quattro, ciascuna con
la propria scadenza:

\begin{center}
\small
\begin{adjustbox}{max width=\textwidth}
\begin{tabular}{@{}lll@{}}
\toprule
\textbf{Stato} & \textbf{Significato} & \textbf{Scadenza} \\
\midrule
\emph{pending} & rivendicata da un chiamante che sta creando l'esecuzione &
  \code{maxLifetime} \\
\emph{published} & associata a un'esecuzione ancora viva & \code{maxLifetime} \\
\emph{abandoned} & pubblicata e poi esplicitamente abbandonata: l'esecuzione
  non girer\`a mai, il legame pu\`o essere rivendicato di nuovo & --- \\
\emph{terminal} & l'esecuzione \`e archiviata: \`e il tombstone & \code{ttl} da
  completamento \\
\bottomrule
\end{tabular}
\end{adjustbox}
\end{center}

Solo lo stato \emph{abandoned} \`e rivendicabile: l'assenza di un record vivo o di
un esito non \`e mai la prova che un legame pubblicato sia stato abbandonato.
La scadenza per stato non \`e configurata a parte, ma derivata dalle proprieta'
dello store delle esecuzioni, e la ragione \`e nel commento: una chiave \`e utile
solo finch\'e esiste la risposta che nomina, e una chiave che scade prima del
suo esito esegue silenziosamente due volte la funzione, che \`e esattamente il
guasto che la chiave esiste per evitare. Il passaggio a \emph{terminal} riavvia
l'orologio da completamento.

\begin{lstlisting}[language=Java, caption={Il ciclo di rivendicazione (schema).}]
public AcquireResult acquireOrGet(String functionName, String key) {
    ...
    // riserva prima uno slot di budget, poi tenta putIfAbsent;
    // se perde la corsa, restituisce la riserva
    ... claimed(token) | pending() | existing(id, terminal) | budgetExhausted()
}
\end{lstlisting}

Gli esiti possibili --- \code{CLAIMED}, \code{PENDING}, \code{EXISTING},
\code{BUDGET\_EXHAUSTED}, \code{MISSING} --- corrispondono a reazioni distinte del
chiamante: procedere alla creazione, attendere, riusare l'esecuzione esistente,
rifiutare, ricominciare. La pubblicazione e l'abbandono sono condizionati al
possesso del token.

Il numero di chiavi \`e limitato da \code{max-keys}, e il budget \`e una
\emph{riserva atomica}, non la lettura di una dimensione: una nuova
rivendicazione riserva il proprio slot \emph{prima} di creare il legame, cos\`i che
chiavi distinte non possano osservare tutte spazio disponibile e inserirsi tutte
(\emph{finding R7}). Ogni voce possiede esattamente uno slot, legato alla
sua identit\`a e non al suo stato, e lo restituisce una volta sola, quando la voce
lascia la cache. Le repliche di una chiave gi\`a posseduta non toccano mai il budget,
quindi la saturazione non le blocca, e il tetto non espelle mai una chiave viva
per far posto a una nuova. Un ingresso rifiutato per budget esaurito \`e contato
(\code{idempotency\_key\_budget\_rejections}) e l'occupazione \`e esposta come
\code{idempotency\_keys\_held}, senza etichette che rivelerebbero un identificatore
di esecuzione o una chiave dell'utente.

Il ciclo di rivendicazione pu\`o attendere sotto contesa, ed \`e questo a motivare il
confinamento sullo scheduler elastico descritto nel
\cref{cap:04-architettura-control-plane}: una richiesta \emph{senza} chiave non lo
attraversa e resta sull'event loop.

\section{Dispatch}

\subsection{Il router e il trasporto}

\code{DispatcherRouter} \`e volutamente banale: due metodi che delegano a
\code{LocalDispatcher} o a \code{ExternalDispatcher}. La modalit\`a
\code{DEPLOYMENT} non ha un dispatcher proprio --- una volta che il deployment
esiste, invocarlo \`e la stessa operazione che invocare un endpoint esterno, e
la differenza sta nel \emph{chi} ha creato l'endpoint
(\cref{cap:07-modalita-esecuzione}).

La scelta della modalit\`a non vive pi\`u nel gestore di completamento:
\code{AttemptTransportAdapter}, nel control plane, \`e l'unico posto da cui il
percorso del tentativo chiama il router e la condizione di prontezza. Il
coordinatore dei tentativi, in \code{execution-runtime}, non importa n\'e l'uno n\'e
l'altra, e il commento spiega perch\'e: il cancello del motore di scheduling aveva
gi\`a dovuto essere indebolito una volta per prendere un lock per funzione, e una
chiamata al provider raggiungibile dalle vicinanze del cancello sarebbe una
classe di problema ben peggiore.

\subsection{\code{ExternalDispatcher} e il contratto sul filo}

Il dispatcher esterno effettua una \code{POST} verso l'endpoint della funzione,
allegando quattro header di contesto: identificatore di esecuzione, numero di
tentativo, trace id e chiave di idempotenza. Il timeout \`e applicato in due
punti: come \code{responseTimeout} sulla richiesta Reactor Netty sottostante,
che interrompe effettivamente la connessione anzich\'e limitarsi ad abbandonare
l'attesa, e come timeout del \code{Mono}, che produce l'errore
\code{EXTERNAL\_TIMEOUT}. Il secondo ha un effetto che il commento del
\code{HttpClientConfig} nomina: cancella anche l'acquisizione pendente dal pool.

Nella direzione di ritorno, il dispatcher interpreta tre header prodotti dal
runtime:

\begin{center}
\small
\begin{adjustbox}{max width=\textwidth}
\begin{tabular}{@{}ll@{}}
\toprule
\textbf{Header} & \textbf{Significato} \\
\midrule
\code{X-Cold-Start} & La richiesta ha pagato un'inizializzazione. \\
\code{X-Init-Duration-Ms} & Durata di tale inizializzazione. \\
\code{X-NanoFaaS-Function-Status} & Il codice di stato HTTP \`e deciso
dall'handler. \\
\bottomrule
\end{tabular}
\end{adjustbox}
\end{center}

Il terzo \`e il pi\`u delicato, e il codice lo tratta con diffidenza esplicita:

\begin{lstlisting}[language=Java]
// An out-of-range status on a marker-bearing response is not spec-legal
// ([200,599] only). Fall through to the platform-error path rather than
// trusting it: EXTERNAL/DEPLOYMENT endpoints are arbitrary unauthenticated
// URLs, and HTTP's status-line grammar is 3DIGIT, so a misbehaving upstream
// can emit one.
\end{lstlisting}

L'endpoint di una funzione \`e un URL arbitrario e non autenticato: il control
plane non pu\`o assumere che rispetti il contratto. Un marcatore che dichiara
``il mio codice di stato \`e intenzionale'' accompagnato da un codice fuori
intervallo viene trattato come errore di piattaforma, non propagato.

\subsection{Il client HTTP: pool, DNS e connessioni}

Il client dei dispatcher non usa il pool globale di Reactor Netty: il pool
(\code{DispatchConnectionPool}) appartiene al contesto e ha un bilancio
dichiarato, per \emph{destinazione}: 500 connessioni, coda di acquisizione pari
a due volte tale valore, timeout di acquisizione 45~s, chiusura delle connessioni
inattive dopo 30~s, e smaltimento dei pool di destinazioni inattive. Il commento
riporta una scelta di \emph{non} adottare un timeout di acquisizione pi\`u breve:
misurato come grande vantaggio solo in un harness che non cancellava mai, e
senza effetto misurabile quando, come fa \code{ExternalDispatcher}, si cancella.

Il client risolve il DNS a \emph{ogni} nuova connessione invece di conservarne le
risposte (\code{NoopDnsCache}; non un TTL zero, che Netty rifiuta). La
motivazione, nel commento, \`e un incidente: una funzione cancellata e
ri-registrata con lo stesso nome riceve un nuovo Service con un nuovo ClusterIP;
con la cache del record (30~s su CoreDNS di minikube) i dispatch continuavano a
chiamare il vecchio ClusterIP, dove nulla rispondeva, e ogni connessione restava
appesa fino al timeout di connessione pur essendo il nuovo pod Ready. Poich\'e le
connessioni sono in pool, la risoluzione avviene solo all'apertura di una nuova.

\section{Riprova}
\label{sec:06-riprova}

La riprova \`e gestita da \code{AttemptCoordinator}, in \code{execution-runtime};
\code{ExecutionCompletionHandler} ne \`e oggi una facciata Spring. La regola di
base non \`e cambiata --- si riprova finch\'e il numero di tentativo non supera
\code{maxRetries} --- ma i dettagli s\`i.

\paragraph{La riprova ha un ritardo.} Nella prima versione una riprova era
immediata, e il suo costo era misurabile: dopo un deploy, gli errori
\code{EXTERNAL\_ERROR} erano riprove ravvicinate contro un endpoint che non era ancora
pronto, che si esaurivano in pochi millisecondi. Oggi \code{RetryBackoff} calcola
l'istante di eleggibilit\`a del prossimo tentativo: la base cresce come
$\min(\text{max},\,\text{iniziale}\cdot 2^{n-1})$ e il ritardo effettivo \`e
campionato con \emph{jitter uguale} in $[\text{base}/2,\,\text{base}]$; il minimo
non \`e mai zero. Un \code{Retry-After} del server a monte (accettato in secondi o
come data, solo per \code{429} e \code{503}, e solo se coerente fra pi\`u valori) \`e un
\emph{minimo}, non un tetto: se \`e pi\`u lontano del ritardo locale, vince.
I default sono \code{nanofaas.retry.initial-backoff}=100~ms e
\code{max-backoff}=2~s.

Una riprova in attesa \emph{conserva la propria riserva} e non occupa capacit\`a di
dispatch: nel profilo diretto il lavoro in attesa \`e tenuto da
\code{ExecutorBackedInvocationEnqueuer}, con un timer e un numero massimo di lavori
pendenti; con un modulo di coda, il biglietto rientra nel motore con \code{notBefore}
e attende nell'insieme dei \emph{ritardati} (\cref{cap:09-async-queue}).

\paragraph{Misure di riprova.} Il confronto \`e in
\code{docs/experiments/scheduler-switching-2026-09} (base \code{a17289c2}). Nella
riproduzione dal vivo di un deploy, la base ha completato 742 esecuzioni su 900 e il
candidato 900 su 900. Con un \code{429} controllato e \code{Retry-After: 1}, la base
ha tentato 1\,195 volte contro 600, ha concluso 3 esecuzioni con successo e 297 con
errore contro 300 e 0, con un intervallo minimo di 0,000727~s contro 1,000698~s e 892
intervalli sotto il secondo contro nessuno. Va detto che un caso su dodici del
deploy vivo resta senza spiegazione: un timeout che i log non chiariscono.

\paragraph{Il cancello di prontezza per DEPLOYMENT.} Un Deployment appena creato
non \`e pronto per i secondi che il primo pod impiega ad avviarsi, e un dispatch
inviato allora incontra un Service senza endpoint. Invece di spendere riprove,
\code{DeploymentWakeUpGate} trattiene il dispatch finch\'e la generazione non \`e stata
osservata pronta, una volta sola: la funzione gi\`a vista pronta costa una ricerca
e nessuna chiamata al provider. L'attesa \`e cancellabile, ha un timeout
(\code{nanofaas.deployment.wakeup.timeout}, 30~s), un intervallo di polling
(250~ms) e un'et\`a massima dell'osservazione di pronto (5~s); la stessa attesa
serve la scale-to-zero.

\paragraph{Tre dettagli che restano.}
Primo, il tentativo di riprova viene costruito come un \code{InvocationTask}
nuovo, con la chiave di idempotenza \emph{azzerata}:

\begin{lstlisting}[language=Java]
InvocationTask retryTask = new InvocationTask(
        ..., null,  // No idempotency key for retry - retry is internal
        ..., currentTask.attempt() + 1, currentTask.kind());
\end{lstlisting}

La riprova \`e un'operazione interna sulla \emph{stessa} esecuzione: conservare
la chiave farebbe s\`i che la riprova si riconosca come duplicato di se stessa e
riproduca l'esito fallito. (Per questo il record ha catturato la chiave alla
costruzione.)

Secondo, il rilascio dello slot di concorrenza \`e idempotente per tentativo. Nella
prima versione era un insieme di tentativi rilasciati nel record; oggi \`e
\code{DispatchLease.release()}, che rilascia una volta sola. Senza questa
protezione, un esito che arriva due volte --- risposta HTTP \emph{e} callback ---
rilascerebbe due slot avendone occupato uno, e il contatore delle richieste in
volo scenderebbe sotto il valore reale, con effetti diretti sull'autoscaler e sul
controllore di concorrenza che leggono quel contatore.

Terzo, la riprova pu\`o fallire perch\'e non c'\`e modo di ripubblicarla (coda piena,
lavori pendenti al massimo, timer rifiutato), e in tal caso l'esecuzione viene
completata con l'errore \emph{originale} anzich\'e con un errore di coda: il
chiamante riceve la ragione per cui la funzione \`e fallita, non la ragione per cui
la piattaforma ha rinunciato a riprovare. Il codice riverifica che l'esecuzione sia
ancora parcheggiata su quel tentativo prima di concluderla, perch\'e una scadenza
amministrativa pu\`o averla conclusa nel frattempo.

\section{Limitazione di frequenza}

Il limitatore \`e ancora una finestra fissa di un secondo, ma ha cambiato sia
posizione sia realizzazione. La posizione: non \`e pi\`u chiamato da
\code{InvocationService}, dopo la decodifica HTTP e JSON, ma da un \code{WebFilter}
(\code{RateLimitWebFilter}) ristretto ai due suffissi \code{:invoke} e \code{:enqueue}, che
rifiuta \emph{prima} di leggere il corpo: nel commento, ogni rifiuto pagava prima
tutto il costo di decodifica, proprio quando la piattaforma ha meno da spendere.
Il rifiuto svuota comunque il corpo, altrimenti Reactor Netty non riceve il segnale
di richiesta consumata e pu\`o chiudere la connessione invece di riusarla in
keep-alive, un costo maggiore di ci\`o che il rifiuto risparmia; risponde \code{429}
con \code{Retry-After: 1}. La realizzazione: finestra e conteggio stanno in un \emph{unico} \code{AtomicLong}
(33 bit per il secondo epoch, 31 per il conteggio), cos\`i che il passaggio di
finestra e il reset del contatore siano una sola transizione CAS. La versione
precedente, con due contatori atomici separati, aveva una corsa che il commento
descrive: rifiutava a torto richieste della nuova finestra e permetteva a un reset
di cancellare incrementi fatti dopo il passaggio.

\begin{lstlisting}[language=Java]
public boolean allow() {
    long max = maxPerSecond;
    long now = epochSecondSource.getAsLong();     // letto una volta per decisione
    while (true) {
        long state = windowAndCount.get();
        long window = state >>> COUNT_BITS, count = state & COUNT_MASK;
        if (now > window) {                       // nuova finestra: la reclama a 1
            if (windowAndCount.compareAndSet(state, encode(now, 1L))) return 1L <= max;
            continue;
        }
        if (count >= max) return false;
        if (windowAndCount.compareAndSet(state, state + 1L)) return true;
    }
}
\end{lstlisting}

Il difetto noto della finestra fissa non \`e cambiato, e il commento della classe
lo dichiara: a cavallo del secondo si pu\`o ammettere fino al doppio del limite, e
non \`e un token bucket. La scelta \`e coerente con il ruolo del limitatore in
questo sistema --- una protezione presente sul percorso ma neutra per default
(\cref{cap:03-requisiti-principi}) --- e la sua imprecisione \`e irrilevante a un
milione di richieste al secondo. Il limite \`e \code{volatile} e riletto a ogni
decisione: l'endpoint di configurazione a caldo lo modifica con effetto alla
decisione successiva, senza sincronizzazione aggiuntiva.

\section{Metriche}

\code{Metrics} incapsula la registrazione Micrometer. La struttura ha una
propriet\`a che sorprende alla prima lettura: i contatori per funzione sono
creati alla registrazione della funzione e \emph{rimossi} alla sua
cancellazione, e ogni metodo di incremento verifica prima che i contatori
esistano ancora.

\begin{lstlisting}[language=Java]
public void success(String function) {
    FunctionMeters metersFor = metersOrNull(function);
    if (metersFor != null) { metersFor.success().increment(); }
}
\end{lstlisting}

Il percorso comune --- una funzione registrata e viva --- non prende lock. Lo
prendeva sempre, ed era un monitor \emph{globale} condiviso da tutte le
funzioni, su un percorso che ogni invocazione attraversa sei volte: il
commento riporta 223~ns per operazione con un thread e 2\,638~ns con otto (il
costo per operazione cresceva con i thread invece di restare piatto, che \`e la
firma di una serializzazione). Il lock resta sul percorso lento, dove serve:
prima registrazione e corsa con la rimozione. Un incremento che afferra i
contatori un istante prima della rimozione va perso, ed \`e un prezzo accettato.

La rimozione dei contatori alla cancellazione evita che il registro Prometheus
cresca indefinitamente in un sistema in cui le funzioni sono create e distrutte da
script di test. Ha per\`o una conseguenza sperimentale che il progetto ha
imparato per via empirica: una verifica che interroghi le metriche \emph{dopo}
aver cancellato la funzione non trova nulla. Il \cref{cap:22-metodo-sperimentale}
riporta come gli scenari di validazione sono stati corretti di conseguenza.
Le metriche di carico per funzione dei moduli (\code{function\_queue\_depth} e
compagne) seguono la stessa regola, ma con \emph{drenaggio}: i contatori vengono
rimossi solo dopo che il lavoro in volo della generazione \`e tornato
(\cref{cap:09-async-queue}).

Le esecuzioni relative a funzioni gi\`a rimosse vengono contabilizzate su un
registro separato e non esportato (\code{removed\_function\_*}), cos\`i che la
loro esistenza non produca n\'e serie orfane n\'e eccezioni. Il conteggio dei
campioni \`e legato alla generazione: un completamento la cui funzione \`e stata
rimossa e ri-registrata dopo l'ammissione non attribuisce il proprio campione ai
contatori della nuova registrazione (\code{isCurrentGeneration}, I7).
Tre contatori per percorso (\code{admitted}, \code{refused}, \code{replayed})
distinguono la porta da cui l'invocazione \`e entrata, e sono nomi \emph{nuovi} e non
un'etichetta sui contatori esistenti: cinque di quelli alimentano anelli di
controllo (la latenza il governatore di concorrenza, i dispatch l'autoscaler, in
volo e profondit\`a l'HPA di Kubernetes), e dividere una serie che un anello legge
cambia ci\`o che l'anello vede.
